Pendant le même temps, les troupes royales quittent Albi et se rendent
à Teillet pour en raser les fortifications. Chemin faisant, 1200
partisans huguenots venus de Castres et Réalmont essaient de leur
barrer la route. Ils se font battre. Et le maréchal de Thémines va
faire le siège de Castres. Alban est requis de fournir sa quote-part
pour subvenir à l'entretien de l'armée du roi. Ici, l'on est pauvre
et l'on se refuse à toute réquisition. Heureusement, une trêve est
signée en 1625. Les consuls d'Alban demandent la réunion d'une
conférence pour en fixer les modalités et préciser les conditions
de mise en application.

Notre compatriote Jehan Pujol de la Tapie, en cette même période,
seconde Condé au siège de Pamiers, puis revient chez lui dans sa
montagne. Montmorency, le gouverneur du Languedoc, va l'y utiliser
comme chef de partisans. Il « pétarde » le château de Lacan, entre
Castres et Roquecourbe, fait sauter « le castellas » de Rayssac et
rançonne le seigneur de Calmels. Il lui extorque 1500 livres :
une bonne prise !!!

Tandis que Jean Pujol de la Tapie « pétarde », Condé s'empare de
Réalmont et vient livrer des escarmouches dans la région de Saint
Affrique. En juin 1628, il traverse Alban et ramène ses troupes à
Albi. Comme bien l'on pense, le soldat se nourrit et se loge comme
il peut sans trop de ménagements. Aussi les habitants d'Alban
veulent-ils se faire indemniser pour les dégâts causés par la
soldatesque du roi. On ne sait point s'ils y parvinrent.